\chapter*{Tabla de seguimiento del alumnado}
\addcontentsline{toc}{chapter}{Tabla de seguimiento del alumnado}

Usa esta tabla para registrar en que capitulos necesitas volver a practicar y en cuales ya
resuelves con autonomia.

\begin{tabularx}{\textwidth}{>{\raggedright\arraybackslash}p{0.42\textwidth} c c >{\raggedright\arraybackslash}X}
\toprule
Bloque & Fecha & Nivel actual & Observacion breve \\
\midrule
\texttt{C01} Numeros reales y herramientas numericas & & & \\
\texttt{C02} Polinomios y fracciones algebraicas & & & \\
\texttt{C03} Ecuaciones y sistemas & & & \\
\texttt{C04} Inecuaciones & & & \\
\texttt{C05} Trigonometria & & & \\
\texttt{C06} Vectores & & & \\
\texttt{C07} Geometria analitica & & & \\
\texttt{C08} Limites, continuidad y asintotas & & & \\
\texttt{C09} Derivadas y aplicaciones & & & \\
Repaso acumulativo y simulacros & & & \\
\bottomrule
\end{tabularx}

\vspace{1em}

\textbf{Clave sugerida}
\begin{itemize}[leftmargin=*]
  \item \texttt{A}: lo hago con autonomia y apenas consulto ayudas.
  \item \texttt{G}: necesito apoyo puntual o revisar un ejemplo.
  \item \texttt{R}: debo repasar la teoria minima y rehacer guiados.
\end{itemize}
